% Options for packages loaded elsewhere
\PassOptionsToPackage{unicode}{hyperref}
\PassOptionsToPackage{hyphens}{url}
\documentclass[
]{article}
\usepackage{xcolor}
\usepackage[margin=1in]{geometry}
\usepackage{amsmath,amssymb}
\setcounter{secnumdepth}{5}
\usepackage{iftex}
\ifPDFTeX
  \usepackage[T1]{fontenc}
  \usepackage[utf8]{inputenc}
  \usepackage{textcomp} % provide euro and other symbols
\else % if luatex or xetex
  \usepackage{unicode-math} % this also loads fontspec
  \defaultfontfeatures{Scale=MatchLowercase}
  \defaultfontfeatures[\rmfamily]{Ligatures=TeX,Scale=1}
\fi
\usepackage{lmodern}
\ifPDFTeX\else
  % xetex/luatex font selection
\fi
% Use upquote if available, for straight quotes in verbatim environments
\IfFileExists{upquote.sty}{\usepackage{upquote}}{}
\IfFileExists{microtype.sty}{% use microtype if available
  \usepackage[]{microtype}
  \UseMicrotypeSet[protrusion]{basicmath} % disable protrusion for tt fonts
}{}
\makeatletter
\@ifundefined{KOMAClassName}{% if non-KOMA class
  \IfFileExists{parskip.sty}{%
    \usepackage{parskip}
  }{% else
    \setlength{\parindent}{0pt}
    \setlength{\parskip}{6pt plus 2pt minus 1pt}}
}{% if KOMA class
  \KOMAoptions{parskip=half}}
\makeatother
\usepackage{color}
\usepackage{fancyvrb}
\newcommand{\VerbBar}{|}
\newcommand{\VERB}{\Verb[commandchars=\\\{\}]}
\DefineVerbatimEnvironment{Highlighting}{Verbatim}{commandchars=\\\{\}}
% Add ',fontsize=\small' for more characters per line
\usepackage{framed}
\definecolor{shadecolor}{RGB}{248,248,248}
\newenvironment{Shaded}{\begin{snugshade}}{\end{snugshade}}
\newcommand{\AlertTok}[1]{\textcolor[rgb]{0.94,0.16,0.16}{#1}}
\newcommand{\AnnotationTok}[1]{\textcolor[rgb]{0.56,0.35,0.01}{\textbf{\textit{#1}}}}
\newcommand{\AttributeTok}[1]{\textcolor[rgb]{0.13,0.29,0.53}{#1}}
\newcommand{\BaseNTok}[1]{\textcolor[rgb]{0.00,0.00,0.81}{#1}}
\newcommand{\BuiltInTok}[1]{#1}
\newcommand{\CharTok}[1]{\textcolor[rgb]{0.31,0.60,0.02}{#1}}
\newcommand{\CommentTok}[1]{\textcolor[rgb]{0.56,0.35,0.01}{\textit{#1}}}
\newcommand{\CommentVarTok}[1]{\textcolor[rgb]{0.56,0.35,0.01}{\textbf{\textit{#1}}}}
\newcommand{\ConstantTok}[1]{\textcolor[rgb]{0.56,0.35,0.01}{#1}}
\newcommand{\ControlFlowTok}[1]{\textcolor[rgb]{0.13,0.29,0.53}{\textbf{#1}}}
\newcommand{\DataTypeTok}[1]{\textcolor[rgb]{0.13,0.29,0.53}{#1}}
\newcommand{\DecValTok}[1]{\textcolor[rgb]{0.00,0.00,0.81}{#1}}
\newcommand{\DocumentationTok}[1]{\textcolor[rgb]{0.56,0.35,0.01}{\textbf{\textit{#1}}}}
\newcommand{\ErrorTok}[1]{\textcolor[rgb]{0.64,0.00,0.00}{\textbf{#1}}}
\newcommand{\ExtensionTok}[1]{#1}
\newcommand{\FloatTok}[1]{\textcolor[rgb]{0.00,0.00,0.81}{#1}}
\newcommand{\FunctionTok}[1]{\textcolor[rgb]{0.13,0.29,0.53}{\textbf{#1}}}
\newcommand{\ImportTok}[1]{#1}
\newcommand{\InformationTok}[1]{\textcolor[rgb]{0.56,0.35,0.01}{\textbf{\textit{#1}}}}
\newcommand{\KeywordTok}[1]{\textcolor[rgb]{0.13,0.29,0.53}{\textbf{#1}}}
\newcommand{\NormalTok}[1]{#1}
\newcommand{\OperatorTok}[1]{\textcolor[rgb]{0.81,0.36,0.00}{\textbf{#1}}}
\newcommand{\OtherTok}[1]{\textcolor[rgb]{0.56,0.35,0.01}{#1}}
\newcommand{\PreprocessorTok}[1]{\textcolor[rgb]{0.56,0.35,0.01}{\textit{#1}}}
\newcommand{\RegionMarkerTok}[1]{#1}
\newcommand{\SpecialCharTok}[1]{\textcolor[rgb]{0.81,0.36,0.00}{\textbf{#1}}}
\newcommand{\SpecialStringTok}[1]{\textcolor[rgb]{0.31,0.60,0.02}{#1}}
\newcommand{\StringTok}[1]{\textcolor[rgb]{0.31,0.60,0.02}{#1}}
\newcommand{\VariableTok}[1]{\textcolor[rgb]{0.00,0.00,0.00}{#1}}
\newcommand{\VerbatimStringTok}[1]{\textcolor[rgb]{0.31,0.60,0.02}{#1}}
\newcommand{\WarningTok}[1]{\textcolor[rgb]{0.56,0.35,0.01}{\textbf{\textit{#1}}}}
\usepackage{longtable,booktabs,array}
\newcounter{none} % for unnumbered tables
\usepackage{calc} % for calculating minipage widths
% Correct order of tables after \paragraph or \subparagraph
\usepackage{etoolbox}
\makeatletter
\patchcmd\longtable{\par}{\if@noskipsec\mbox{}\fi\par}{}{}
\makeatother
% Allow footnotes in longtable head/foot
\IfFileExists{footnotehyper.sty}{\usepackage{footnotehyper}}{\usepackage{footnote}}
\makesavenoteenv{longtable}
\usepackage{graphicx}
\makeatletter
\newsavebox\pandoc@box
\newcommand*\pandocbounded[1]{% scales image to fit in text height/width
  \sbox\pandoc@box{#1}%
  \Gscale@div\@tempa{\textheight}{\dimexpr\ht\pandoc@box+\dp\pandoc@box\relax}%
  \Gscale@div\@tempb{\linewidth}{\wd\pandoc@box}%
  \ifdim\@tempb\p@<\@tempa\p@\let\@tempa\@tempb\fi% select the smaller of both
  \ifdim\@tempa\p@<\p@\scalebox{\@tempa}{\usebox\pandoc@box}%
  \else\usebox{\pandoc@box}%
  \fi%
}
% Set default figure placement to htbp
\def\fps@figure{htbp}
\makeatother
\setlength{\emergencystretch}{3em} % prevent overfull lines
\providecommand{\tightlist}{%
  \setlength{\itemsep}{0pt}\setlength{\parskip}{0pt}}
\setlength{\emergencystretch}{4em}
\usepackage{etoolbox}
\makeatletter
\renewcommand{\maketitle}{\begin{titlepage}\centering\vspace*{1.1in}{\LARGE\bfseries\@title\par}\vspace{2.25em}{\large\@author\par}\vspace{1em}{\normalsize\textit{Faculty Advisor: Dr. Ali Duman}\par}\vspace{0.6em}{\normalsize\textit{Committee Members: Dr. Katherine Shoemaker, Dr. Timothy Redl}\par}\vfill\end{titlepage}\clearpage}
\makeatother
\apptocmd{\tableofcontents}{\clearpage}{}{}
\AtBeginEnvironment{longtable}{\footnotesize\setlength{\tabcolsep}{4pt}}
\setlength{\hfuzz}{6pt}
\IfFileExists{fvextra.sty}{\usepackage{fvextra}\RecustomVerbatimEnvironment{Highlighting}{Verbatim}{commandchars=\\\{\},breaklines=true,breakanywhere=true}}{}
\AtBeginDocument{\sloppy}
\usepackage{bookmark}
\IfFileExists{xurl.sty}{\usepackage{xurl}}{} % add URL line breaks if available
\urlstyle{same}
\hypersetup{
  pdftitle={Graph-Based Screening for Potentially Cooperative Vessel Pairs},
  pdfauthor={Oumayma Ben Aoun},
  hidelinks,
  pdfcreator={LaTeX via pandoc}}

\title{Graph-Based Screening for Potentially Cooperative Vessel Pairs}
\usepackage{etoolbox}
\makeatletter
\providecommand{\subtitle}[1]{% add subtitle to \maketitle
  \apptocmd{\@title}{\par {\large #1 \par}}{}{}
}
\makeatother
\subtitle{A temporal link prediction approach using AIS (2012--2019)}
\author{Oumayma Ben Aoun}
\date{}

\begin{document}
\maketitle

{
\setcounter{tocdepth}{3}
\tableofcontents
}
\section*{Abstract}\label{abstract}
\addcontentsline{toc}{section}{Abstract}

Illegal, unreported, and unregulated (IUU) fishing has increased
interest in tools that surface vessels that may operate in a coordinated
way. This project develops an \textbf{unsupervised} screening pipeline
built on Automatic Identification System (AIS) data (2012--2019). Each
day is represented as a graph of vessels; a \textbf{Temporal Graph
Convolutional Network (TGCN)} performs \textbf{temporal link prediction}
so that pairs that plausibly co-occur in the future rank highly. Raw
tracks are then checked independently with simple geographic rules to
prioritize pairs for analyst review. \textbf{Training graph edges}
connect vessels whose positions fall within \textbf{10 km} on the
\textbf{same calendar day}. Wider distance checks (for example
\textbf{25 km}) are used only \textbf{after} scoring to validate
candidates---they are \textbf{not} training labels. The model does
\textbf{not} receive labels for ``cooperative'' behavior; it learns from
\textbf{structure and timing} in the graphs and produces
\textbf{candidates for review}, not adjudicated outcomes.

\textbf{Main quantitative result (Table 1):} On the \textbf{combined}
graph---proximity edges plus optional \textbf{social} edges linking
vessels that share the same three-digit maritime identifier (MID)
prefix, implemented with a corrected MID rule---a reproducible laptop
configuration (\textbf{1450} training days, 3 epochs, seed 1) achieves
\textbf{ROC AUC = 0.72997} and \textbf{average precision (AP) = 0.71381}
on \textbf{876} held-out daily test buckets
(\texttt{artifacts/tgcn\_social\_maxb1450\_ep3.json}). A
\textbf{proximity-only} ablation on the same edge timeline reaches
higher scores (about \textbf{0.78} AUC) but on a much smaller
\textbf{20-bucket evaluation slice}; it is reported in Table 1 as a
\textbf{directional ablation}, not a like-for-like headline. A narrow
\textbf{August 2017} window was used early for exploration; its metrics
\textbf{cannot} be compared with the full 2012--2019 benchmark, and this
report does not include figures from that window. After strict
\textbf{25 km} and \textbf{±1 day} screening, \textbf{eight} pairs
remain of interest; a \textbf{gear-aware} follow-up analysis suggests
that many \textbf{excluded} top candidates resemble routine same-fleet
trawler co-location rather than unusual encounters. We \textbf{do not}
assert legal conclusions or membership in any ``dark fleet.''

\textbf{Takeaways.} \textbf{(i)} Modest absolute scores are
\textbf{expected} for an unsupervised, label-free screening task on
global AIS, and the ±0.10 day-to-day spread is a property of the domain
rather than a flaw of the model. \textbf{(ii)} \textbf{Post-hoc}
\textbf{25 km / ±1 day} geography matters because \textbf{10 km}
same-day graph structure can elevate fleet-like pairs that rarely meet
the tighter physical standard; the geographic filter separates
\textbf{fleet-scale similarity} from \textbf{pair-level proximity}.
\textbf{(iii)} \textbf{How} registry information enters the graph
matters as much as \textbf{whether} it does: dense same-MID social edges
add noise, encourage over-smoothing, and compete with proximity in the
loss---future work should fuse registry context as \textbf{node
features} or as \textbf{typed edges} instead. \textbf{(iv)} The
deliverable is a \textbf{short, auditable} list (top-200 → 8 geography
survivors, \textasciitilde4\%; as a group, cross-gear pairs survive at
\textasciitilde5.1\% vs 2.4\% for trawler--trawler, and the best
cross-gear category at \textasciitilde10× the trawler--trawler rate)
with \textbf{tracks and overlap plots}---\textbf{prioritized leads} plus
\textbf{supporting context}, \textbf{not} a determination of misconduct.

\textbf{Keywords:} AIS, IUU fishing, temporal graph neural networks,
link prediction, vessel co-movement

\begin{center}\rule{0.5\linewidth}{0.5pt}\end{center}

\section{Introduction}\label{introduction}

\textbf{Problem.} Fisheries enforcement and marine research often need
to narrow attention to vessel pairs that may interact at sea---whether
through transshipment, repeated joint presence, or other coordinated
patterns sometimes associated with periods of missing AIS (``dark''
behavior). AIS offers broad spatial and temporal coverage, yet
\textbf{explicit labels for cooperation are seldom available} in
operational datasets.

\textbf{What this work does.} We cast screening as \textbf{temporal link
prediction} on \textbf{daily} graphs of vessels. Pairs that receive high
scores under this objective are \textbf{candidates} for closer
inspection. We support that step with straightforward \textbf{post-hoc}
checks: distances between \textbf{daily mean positions}, maps of raw
tracks, and summaries of overlap over time.

\textbf{What this work does not do.} This pipeline does \textbf{not}
establish coordination, wrongdoing, or IUU activity. It yields a
\textbf{ranked list} and \textbf{supporting geographic context} intended
for expert review.

\textbf{Contributions (short).}

\begin{enumerate}
\def\labelenumi{\arabic{enumi}.}
\tightlist
\item
  An end-to-end, reproducible workflow: AIS records → daily graphs →
  TGCN with optional temporal node features → chronological train/test
  split.\\
\item
  A \textbf{three-digit MID--consistent} design for combining proximity
  edges with same-prefix \textbf{social} edges
  (\texttt{scripts/mmsi\_mid.py},
  \texttt{scripts/add\_social\_edges.py}).\\
\item
  Reported metrics on \textbf{full-coverage} graphs within
  \textbf{practical laptop memory} limits (training is capped by day
  count where necessary).\\
\item
  Interpretation beyond raw scores: geographic filters, overlap
  visualizations for eight validated pairs, and \textbf{gear-aware
  stratification} (§4.2.1) that compares candidates to fleet-wide
  co-location patterns so that routine fleet behavior can be
  distinguished from unusual pairs.
\end{enumerate}

\textbf{Roadmap.} Section~3 presents \textbf{methods}: formal setup
(\textbf{Appendix C}), \textbf{data}, \textbf{graph construction} (10~km
proximity plus optional MID-based social edges), \textbf{training}, and
\textbf{validation-only} distance rules. Section~4 presents
\textbf{results}: headline metrics (\textbf{Table~1}), an encoder
comparison (\textbf{GCN} vs.~\textbf{graph transformer}) on a
\textbf{smaller capped} graph (§4.1.1; distinct from the full-coverage
experiments in Table~1), geographic and case-study material
(\textbf{Tables~2--3}, Figures~1--4), and gear definitions
(\textbf{Appendix~A}). Section~5 discusses limitations (grouped under
task, data, evaluation, and deployment) and ends with a short
\textbf{synthesis}; Section~6 concludes; \textbf{Appendix~B} collects
reproduction commands.

\begin{quote}
\textbf{For readers and committees (three sentences).} (1)
\textbf{Training} graphs use \textbf{10~km} same-day proximity (plus
optional MID social edges); \textbf{25~km} and wider screens are
\textbf{validation}, not supervised labels. (2) The headline
\textbf{0.72997 / 0.71381} values apply to the \textbf{combined}
full-coverage row in \textbf{Table~1}; \textbf{proximity-only} and
\textbf{capped} settings are \textbf{separate experiments}---see §3.6.
(3) \textbf{Figure~2} uses 2012--2019 fleet context to show that many
high-scoring pairs align with ordinary fleet co-location; the
\textbf{eight} pairs that pass strict geography stand out relative to
that background.
\end{quote}

\begin{center}\rule{0.5\linewidth}{0.5pt}\end{center}

\section{Related work}\label{related-work}

\textbf{Fisheries monitoring and AIS.} Research on illegal, unreported,
and unregulated (IUU) fishing and on coordinated vessel behavior makes
heavy use of \textbf{Automatic Identification System (AIS)} data to
characterize fishing effort and vessel encounters (see References for
global fisheries footprint and transshipment studies). AIS alone does
not encode whether two vessels cooperated or broke rules; analysts
usually combine tracks with \textbf{proximity logic}, \textbf{registry
information}, or \textbf{enforcement intelligence}. That limitation
motivates methods that \textbf{rank pairs without cooperation labels}.

\textbf{AIS-based fisheries and encounter analytics.} Beyond mapping
where vessels go, AIS supports \textbf{encounter analysis} and
\textbf{behavioral screening} (e.g., transshipment and
loitering---Miller et al., 2018; assessments of monitoring
systems---Park et al., 2020). Much of that work uses \textbf{explicit
rules} or \textbf{supervised} signals when labels exist. Here the
learning stage remains \textbf{unsupervised}; geographic checks enter
only as \textbf{post-hoc} validation.

\textbf{Graph learning and link prediction.} Vessel co-presence maps
naturally to a \textbf{graph}: nodes are vessels; edges encode same-day
proximity (and here, optional ties based on registry prefix).
\textbf{Graph convolutional networks} (Kipf \& Welling, 2017) and
\textbf{inductive} neighborhood models (Hamilton et al., 2017) learn
node embeddings from topology and features. \textbf{Link prediction}
ranks missing or future edges (Lü \& Zhou, 2011). When the graph changes
over time, \textbf{temporal} formulations apply---for example
spatio-temporal convolutions (Yu et al., 2018) or recurrent temporal
graph networks (Rossi et al., 2020). Our model follows a
\textbf{TGCN}-style recurrent update over daily snapshots and scores
pairs with an \textbf{inner-product} link head (§3.4).

\textbf{Evaluation philosophy.} Train and test sets split
\textbf{calendar days} in chronological order: the model trains on
earlier periods and is scored on \textbf{later} periods. ROC AUC and
average precision therefore reflect \textbf{forward-looking} ranking
quality. Random mixing of days across time would ignore drift and could
\textbf{inflate} scores; we avoid that design.

\begin{center}\rule{0.5\linewidth}{0.5pt}\end{center}

\section{Methods}\label{methods}

\subsection{Problem formulation}\label{problem-formulation}

Fix a \textbf{calendar day} (time bucket) \(t\). Let \(G_t = (V, E_t)\)
be an \textbf{undirected} graph for that day: \textbf{vertices} \(V\)
are vessel identifiers (\textbf{MMSIs}); \textbf{edges} \(E_t\) are
built as in §3.3---mainly pairs whose positions lie within
\textbf{10~km} great-circle distance on day \(t\), optionally augmented
by \textbf{social} edges between vessels sharing the same
\textbf{three-digit MID} prefix. The learning problem is
\textbf{temporal link prediction}: a \textbf{TGCN} maps the sequence of
graphs to node embeddings; \textbf{pair scores} are inner products of
embeddings. Training minimizes \textbf{binary cross-entropy with logits}
so that \textbf{observed} edges score above \textbf{randomly sampled
non-edges} on each day. There is \textbf{no} label for ``cooperative''
fishing---only structural positives and negatives (see
\texttt{scripts/run\_tgcn\_time\_multiseed.py}; architecture §3.4).
\textbf{Training} and \textbf{test} sets split \textbf{days} in time
order (\textbf{about 70\% / 30\%}). \textbf{ROC AUC} and \textbf{AP}
summarize ranking quality on \textbf{future} test days.

\textbf{Appendix C} collects \textbf{notation}, repeats the formal
picture for readers who prefer a tabular glossary, and distinguishes
\textbf{10~km training edges} from \textbf{validation-only} distances.

\subsection{Data and time buckets}\label{data-and-time-buckets}

AIS records cover \textbf{2012--2019} at \textbf{daily} aggregation (on
the order of \textbf{2.38 billion} vessel--day rows, with sea-surface
temperature merged where available). The study timeline is partitioned
into \textbf{one graph per calendar day}; \textbf{876} days carry the
edge data used for the primary evaluation.

\subsection{Graph construction}\label{graph-construction}

\begin{itemize}
\tightlist
\item
  \textbf{Nodes:} one per \textbf{MMSI} (vessel identifier).\\
\item
  \textbf{Proximity edges:} same \textbf{calendar day},
  \textbf{Haversine} distance at most \textbf{10~km} between
  \textbf{daily} position summaries for the two vessels. The
  full-coverage list \texttt{edges\_2012\_2019\_full.parquet} is
  produced by \texttt{scripts/build\_temporal\_graph\_baseline.py}. Some
  builds \textbf{limit edges per day} to fit in memory; see
  \texttt{README.md} for commands tied to each artifact.\\
\item
  \textbf{Social edges:} undirected ties between vessels that share the
  \textbf{same three-digit ITU MID} prefix.
  \textbf{\texttt{scripts/mmsi\_mid.py}} extracts the prefix so it is
  not truncated at six digits. Social edges are \textbf{capped per day}
  (for example \textbf{2000} per bucket in
  \texttt{scripts/add\_social\_edges.py}) so that graphs remain
  tractable.\\
\item
  \textbf{Combined (proximity + social)} benchmark list:
  \texttt{artifacts/edges\_full\_with\_social.parquet}.
\end{itemize}

\textbf{Validation distances (not training edges).} Candidate review
uses \textbf{25 / 50 / 100~km} thresholds (and sometimes \textbf{±1~day}
alignment) on \textbf{daily mean} positions. These rules \textbf{do not}
define training edges and are \textbf{not} the same as the
\textbf{10~km} edge construction; see §3.5.

\subsection{Model and training
objective}\label{model-and-training-objective}

\textbf{Architecture.} A Temporal Graph Convolutional Network
(\textbf{TGCN}) with optional \textbf{temporal node features} derived
from recent graph statistics (degree, interaction counts, partner
diversity, recency, gaps between events---built in
\texttt{scripts/build\_temporal\_node\_features.py}; summarized in
\texttt{README.md}).

\textbf{Training details (primary laptop benchmark).} Unless stated
otherwise: \textbf{embedding size 32}, \textbf{Adam}, learning rate
\textbf{0.001}, \textbf{3} passes over training days in order,
\textbf{binary cross-entropy with logits} with one \textbf{random
non-edge} sampled per positive edge \textbf{per day}. Chronological
split \textbf{about 70\% train / 30\% test}
(\texttt{-\/-test-ratio\ 0.3}). For the headline \textbf{combined} graph
run, training was limited to \textbf{1450} days because \textbf{larger}
day counts exceeded available \textbf{RAM} on the laptop and led to
out-of-memory termination.

\textbf{Evaluation.} \textbf{ROC AUC} and \textbf{AP} are computed on
\textbf{held-out test days}. The implementation drops or masks edges so
that the test loss does not reuse training positives in a trivial way
(see comments in the TGCN scripts).

\textbf{Computational cost.} Experiments used one \textbf{laptop} (Apple
M-series, \textbf{16~GB} unified memory). The main combined-graph run
(\textbf{1450} train days, \textbf{3} epochs) typically finished in
\textbf{20--40~minutes} depending on system load; \textbf{peak memory}
was about \textbf{12--14~GB}, which explains why caps near \textbf{1480}
days triggered failures. Proximity-only full graphs were trained under
similar limits. Smaller capped graphs and the brief August~2017 run
finished in \textbf{under five minutes} each. All computation was
\textbf{CPU-only} (PyTorch). Timing and memory are \textbf{indicative};
scripts do not log wall-clock or RSS automatically, so exact timings
require a controlled rerun.

\subsection{Heuristic validation (not supervised
labels)}\label{heuristic-validation-not-supervised-labels}

For scored pairs we measure \textbf{great-circle distance} between
\textbf{daily mean positions} and count days falling within \textbf{25 /
50 / 100~km} (and optional \textbf{±1~day} alignment). These steps
describe \textbf{where} vessels were relative to each other; they
\textbf{do not} define the learning objective. §4.2.1 adds
\textbf{gear-aware} context by comparing candidates to
\textbf{fleet-wide} co-location patterns.

\subsection{Which numbers are ``the same
experiment''?}\label{which-numbers-are-the-same-experiment}

Treat each row below as a \textbf{distinct} experimental
setting---\textbf{do not} combine metrics across rows:

{\def\LTcaptype{none} % do not increment counter
\begin{longtable}[]{@{}
  >{\raggedright\arraybackslash}p{(\linewidth - 2\tabcolsep) * \real{0.5714}}
  >{\raggedright\arraybackslash}p{(\linewidth - 2\tabcolsep) * \real{0.4286}}@{}}
\toprule\noalign{}
\begin{minipage}[b]{\linewidth}\raggedright
Setting
\end{minipage} & \begin{minipage}[b]{\linewidth}\raggedright
Role
\end{minipage} \\
\midrule\noalign{}
\endhead
\bottomrule\noalign{}
\endlastfoot
\textbf{2012--2019 combined graph (proximity + social)} &
\textbf{Primary benchmark}---full timeline, social edges, reported
AUC/AP on 876 test buckets. \\
\textbf{Proximity-only graph} & \textbf{Ablation}---no social edges;
often higher AUC on this codebase; shows contribution of the social
layer. \\
\textbf{Capped / smaller graphs} & \textbf{Sanity checks and
tuning}---easier prediction; \textbf{very high} AUC; \textbf{not} the
same difficulty as full global coverage. \\
\textbf{August 2017 short window} & \textbf{Exploratory only}---few
buckets; metrics are \textbf{not} comparable to the main table. No
figures from this window appear in the final report. \\
\end{longtable}
}

\begin{center}\rule{0.5\linewidth}{0.5pt}\end{center}

\section{Results}\label{results}

The presentation follows a standard empirical flow: \textbf{overall
scores} (§4.1), \textbf{maps and interpretation} (§4.2--4.3), and a
brief \textbf{exploratory} slice from August~2017 (§4.4). As emphasized
in §3.6, \textbf{each row of Table~1} corresponds to a
\textbf{different} graph setup---scores from different rows
\textbf{cannot} be merged into one narrative.

\subsection{Quantitative performance
(TGCN)}\label{quantitative-performance-tgcn}

\textbf{Headline numbers, in one place.} On the full \textbf{2012--2019}
combined graph the primary laptop run reports \textbf{ROC AUC = 0.72997}
and \textbf{AP = 0.71381} on \textbf{876} held-out daily test buckets.
The \textbf{proximity-only} ablation on the same edge timeline reaches
\textbf{0.77618 / 0.79549}, but on a much smaller \textbf{20-bucket
evaluation slice} (single-seed, 5 epochs); it is reported as a
\textbf{directional ablation}, not a like-for-like headline. The smaller
\textbf{capped} graph used for tuning reaches roughly \textbf{0.95 /
0.96}, and a \textbf{5-fold rolling} version of that capped setup gives
\textbf{0.927 ± 0.067 / 0.947 ± 0.048} across seeds \textbf{1--5}.
\textbf{Table 1} collects every setting in one block; the rest of this
section explains how to read each row.

\textbf{Table 1.} All TGCN results in one place. Each row is a
\textbf{different} experiment---do not mix metrics across rows (§3.6).
Parquet filenames are omitted here so the table fits the page; see §3
(graph construction) and \textbf{Appendix B} (reproducibility) for full
paths.

{\def\LTcaptype{none} % do not increment counter
\begin{longtable}[]{@{}
  >{\raggedright\arraybackslash}p{(\linewidth - 8\tabcolsep) * \real{0.1698}}
  >{\raggedright\arraybackslash}p{(\linewidth - 8\tabcolsep) * \real{0.3585}}
  >{\raggedright\arraybackslash}p{(\linewidth - 8\tabcolsep) * \real{0.2075}}
  >{\raggedright\arraybackslash}p{(\linewidth - 8\tabcolsep) * \real{0.1698}}
  >{\raggedright\arraybackslash}p{(\linewidth - 8\tabcolsep) * \real{0.0943}}@{}}
\toprule\noalign{}
\begin{minipage}[b]{\linewidth}\raggedright
Setting
\end{minipage} & \begin{minipage}[b]{\linewidth}\raggedright
Train / eval scope
\end{minipage} & \begin{minipage}[b]{\linewidth}\raggedright
Test buckets
\end{minipage} & \begin{minipage}[b]{\linewidth}\raggedright
ROC AUC
\end{minipage} & \begin{minipage}[b]{\linewidth}\raggedright
AP
\end{minipage} \\
\midrule\noalign{}
\endhead
\bottomrule\noalign{}
\endlastfoot
Proximity-only, full coverage edges & full 2012--2017-08 train; eval on
a \textbf{20-bucket} slice (single seed, 5 epochs) & 20 &
\textbf{0.77618} & \textbf{0.79549} \\
\textbf{Combined + social --- primary laptop run} & \textbf{1450} train
buckets, 3 epochs, seed 1 & \textbf{876} & \textbf{0.72997} &
\textbf{0.71381} \\
Capped graph (single seed) & tuning / easier task & --- &
\textasciitilde0.95 & \textasciitilde0.96 \\
Capped graph, 5-fold rolling CV & seeds 1--5 & --- & 0.927 ± 0.067 &
0.947 ± 0.048 \\
August 2017 (case-study window) & exploratory short window & 5 &
\textasciitilde0.56 & \textasciitilde0.67 \\
\end{longtable}
}

\textbf{How to read Table 1.} \textbf{The only headline number is the
combined+social row} (\textbf{0.72997 / 0.71381} on \textbf{876} test
buckets). The proximity-only row sits at \textbf{0.77618 / 0.79549} but
on a \textbf{20-bucket evaluation slice}---roughly \textbf{40× fewer}
test days---so the higher value is \textbf{directional, not
like-for-like}: it says proximity edges score better than
proximity+social on these days, but does not say a fully comparable
proximity-only run would land at 0.78. The \textbf{capped} rows look
much higher because the graph is a \textbf{smaller, denser subset} that
is structurally easier to predict; they are reported only to support
\textbf{tuning} and \textbf{multi-seed sanity checks}, not as
alternative headlines. The \textbf{August 2017} row covers a
\textbf{5-day} window used for early exploration and gear enrichment
(§4.4); it is intentionally \textbf{not} comparable to the main rows.

\textbf{Are these numbers good?} For an \textbf{unsupervised,
label-free} screening task on \textbf{global AIS} with millions of
node-day combinations, \textbf{AUC around 0.73} with \textbf{AP around
0.71} is in line with what should be expected. The \textbf{AP} result is
the more telling figure: on a sparse temporal graph where the vast
majority of candidate pairs are non-edges, an AP near \textbf{0.71}
corresponds to a model that consistently surfaces real co-occurrence
near the top of its rankings---well above the \textbf{prevalence
baseline} that random ranking would yield. The proximity-only
\textbf{directional ablation} (\textasciitilde{}\textbf{0.78} AUC on a
20-bucket evaluation slice---see Table 1 caveat) hints at an upper bound
on what this architecture can extract from purely geographic structure;
closing the remaining gap will likely require \textbf{architectural
changes}, not just longer training (see ``Why do social edges lower
AUC?'' below and §6).

\textbf{Day-to-day variability is a property of the domain, not a bug.}
The per-bucket standard deviations of about \textbf{±0.103} ROC AUC and
\textbf{±0.106} AP across the \textbf{876} test days mean that
\textbf{some days are genuinely easier than others}: graphs are denser
when vessels concentrate on seasonal grounds and sparser during transit
windows; weather, holidays, and policy changes also shift activity. A
\textbf{macro-average} of \textbf{0.73 / 0.71} across that distribution
is the right summary; a stable model on global AIS should not produce
identical performance every single day, and one that did would be a flag
for \textbf{leakage} rather than for skill.

\textbf{Operationally relevant top-N behavior.} For a screening tool the
practical question is: of the pairs the model ranks highest, \textbf{how
many} survive an analyst's independent geographic test? Among the
\textbf{top 200} pairs from the full-coverage combined-graph run,
\textbf{eight} clear the strict \textbf{25 km / ±1 day} rule on raw
daily means---a roughly \textbf{4\%} survival rate on the long list
(Table 2; §4.2)---well above the underlying prevalence of repeated 25 km
contacts among \textbf{all} active pairs, which is orders of magnitude
smaller. The \textbf{gear-aware} stratification in §4.2.1 tightens this
further: as a \textbf{group}, cross-gear combinations survive at
\textbf{5.1\%} versus \textbf{2.4\%} for trawler--trawler (a
\textasciitilde2× aggregate effect), and at the \textbf{extremes} the
best cross-gear category (\textbf{Fixed gear + Trawlers}) survives at
\textbf{25\%}---roughly \textbf{10×} the trawler--trawler rate. The full
funnel is summarized below.

\textbf{Table 1b --- Operational funnel (combined graph, full-coverage
run).} A precision-at-K-style summary of how the top-200 ranked
candidates collapse to an analyst-reviewable list. Numbers come from the
same primary run as Table 1 (\texttt{tgcn\_social\_maxb1450\_ep3.json})
and the gear-aware analysis in §4.2.1.

{\def\LTcaptype{none} % do not increment counter
\begin{longtable}[]{@{}
  >{\raggedright\arraybackslash}p{(\linewidth - 4\tabcolsep) * \real{0.1591}}
  >{\raggedleft\arraybackslash}p{(\linewidth - 4\tabcolsep) * \real{0.1591}}
  >{\raggedright\arraybackslash}p{(\linewidth - 4\tabcolsep) * \real{0.6818}}@{}}
\toprule\noalign{}
\begin{minipage}[b]{\linewidth}\raggedright
Stage
\end{minipage} & \begin{minipage}[b]{\linewidth}\raggedleft
Pairs
\end{minipage} & \begin{minipage}[b]{\linewidth}\raggedright
Pass rate vs.~previous stage
\end{minipage} \\
\midrule\noalign{}
\endhead
\bottomrule\noalign{}
\endlastfoot
Top-200 model-ranked candidates (input list) & 200 & --- \\
Pass strict \textbf{25 km / ±1 day} geographic filter & \textbf{8} &
\textbf{4.0\%} of 200 \\
~~--- within passing set: \textbf{cross-gear} pairs & 6 & \textbf{5.1\%}
of 118 cross-gear candidates \\
~~--- within passing set: \textbf{trawler--trawler} pairs & 2 &
\textbf{2.4\%} of 82 trawler--trawler candidates \\
Highest-discriminating cross-gear category: \textbf{Fixed gear +
Trawlers} & 1 / 4 & \textbf{25.0\%} (\textasciitilde10×
trawler--trawler) \\
\end{longtable}
}

ROC AUC and AP are reported in Table 1 because they are standard, but
the \textbf{top-200 → geography → gear-context} funnel is what makes the
output usable for an analyst.

\textbf{Primary run details (combined row).} Embedding size \textbf{32},
learning rate \textbf{0.001}, temporal node features \textbf{on}.
Primary outputs: \texttt{tgcn\_social\_maxb1450\_ep3.json} plus the
companion \texttt{.csv} in \texttt{artifacts/}. Raising the training day
cap from \textbf{400} toward \textbf{1450} improved metrics until
\textbf{RAM} limited further growth (roughly \textbf{1480--1500} days
caused failure). Other tuning experiments are summarized in
\texttt{tgcn\_improvement\_suite\_summary.csv} (same folder).

\textbf{Proximity-only row (directional ablation).} Values
\textbf{0.77618 / 0.79549} come from
\texttt{tgcn\_time\_temporal\_nodes\_fullcoverage.json} in
\texttt{artifacts/}. That run used the full \textbf{2012--2019} edge
file and trained up to the same train/test cutoff, but
\textbf{evaluated} on only a \textbf{20-bucket evaluation slice} (versus
\textbf{876} buckets for the combined run); reporting it alongside the
combined headline is informative for \textbf{direction} (proximity edges
score higher than proximity+social on these days) but the
\textbf{absolute} numbers are not directly comparable. Re-running the
proximity-only configuration on the full \textbf{876}-bucket horizon is
queued under §6.

\textbf{Variability and uncertainty.} The headline combined result uses
\textbf{one} random seed (seed~\textbf{1}); a full multi-seed interval
for that full-coverage row is \textbf{explicitly noted as future work}
in §6. Two views of spread are already available and are reported as a
partial substitute:

\begin{itemize}
\tightlist
\item
  \textbf{Day-to-day spread (within one seed).} Over the \textbf{876}
  test days, standard deviations of \textbf{per-day} ROC AUC and AP are
  about \textbf{±0.103} and \textbf{±0.106}. Easy and hard days coexist;
  the reported \textbf{0.72997 / 0.71381} is a \textbf{macro-average}
  over days, not a promise about every individual day. The breadth of
  this spread is itself a feature of the domain---seasonal grounds,
  weather, holidays, and policy changes all shift activity---and a model
  that produced identical performance on every test day would be a flag
  for \textbf{leakage} rather than for skill. Per-day CSV:
  \texttt{tgcn\_social\_maxb1450\_ep3.csv} in \texttt{artifacts/}.
\item
  \textbf{Cross-seed spread on a capped graph.} A \textbf{5-fold}
  rolling validation with seeds \textbf{1--5} on the smaller capped
  graph gives \textbf{0.927 ± 0.067} AUC and \textbf{0.947 ± 0.048} AP
  (Table 1). Cross-seed variance there is modest, which is
  \textbf{suggestive but not conclusive} that seed effects on the
  full-coverage row would be smaller than the day-to-day effects above.
  Re-running the headline configuration with \textbf{at least 3--5
  seeds} is the single most useful next experiment for this report and
  is a top item in §6.
\end{itemize}

\subsubsection{GCN vs.~graph
transformers}\label{gcn-vs.-graph-transformers}

Keeping the §3.4 training recipe fixed, \textbf{only} the
\textbf{per-day graph encoder} changes. Training driver:
\texttt{run\_tgcn\_time\_multiseed.py}. Encoder code:
\texttt{temporal\_graph\_baselines.py} in the \texttt{scripts/} tree.
\textbf{two \texttt{GCNConv} layers + GRU} versus \textbf{two
\texttt{TransformerConv} layers + GRU}, with the same
\textbf{dot-product} link head. \textbf{``Graph transformer''} here
means \textbf{neighbor attention} through
\textbf{\texttt{TransformerConv}} at each snapshot---not
\textbf{full-graph} self-attention over all vessels (that would be
\textbf{quadratic} in fleet size and is not attempted).

\textbf{How to read this comparison.}

\begin{enumerate}
\def\labelenumi{\arabic{enumi}.}
\tightlist
\item
  All numbers below use the \textbf{capped} edge list (parquet basename
  contains \texttt{cap5000\_even30}; full path in \textbf{Appendix B}).
  Do \textbf{not} treat this as a second headline for the full-coverage
  \textbf{Table~1} row (§3.6).\\
\item
  The main benchmark uses \textbf{\texttt{-\/-model\ tgcn}}. This
  ablation compares \textbf{\texttt{-\/-model\ gcn}} to
  \textbf{\texttt{-\/-model\ graph\_transformer}},
  i.e.~\textbf{convolution vs.~neighbor attention} in the \textbf{same
  two-layer + GRU} template---not a different temporal core.
\end{enumerate}

\textbf{Protocol.} Train/test day caps \textbf{400} / \textbf{80};
\textbf{5} epochs (the Table~1 combined row uses \textbf{3}); seeds
\textbf{1--3}; embedding \textbf{32}; temporal node features on; hard
negatives; \textbf{49} test days after masking. Write-ups and scripts:
\texttt{gcn\_vs\_graph\_transformer.md} and
\texttt{compare\_gcn\_graph\_transformer.py} (see \textbf{Appendix B}
for full paths).

\emph{Means ± standard deviation are \textbf{across seeds} (1--3),
unlike the day-level spread described for the headline run in §4.1.}

{\def\LTcaptype{none} % do not increment counter
\begin{longtable}[]{@{}lll@{}}
\toprule\noalign{}
Model & ROC AUC & Average precision \\
\midrule\noalign{}
\endhead
\bottomrule\noalign{}
\endlastfoot
\textbf{GCN} + GRU & \textbf{0.6201 ± 0.0102} & \textbf{0.7672 ±
0.0226} \\
\textbf{TransformerConv} + GRU & 0.4576 ± 0.1197 & 0.6735 ± 0.0554 \\
\end{longtable}
}

On average, \textbf{GCN} achieves higher ROC AUC and AP and
\textbf{lower} seed-to-seed variance. Paired tests by calendar day favor
GCN for \textbf{each} seed (see artifact). A plausible---but
\textbf{not} proven---explanation is that \textbf{noisy} proximity
neighborhoods are stabilized more easily by \textbf{fixed} neighborhood
aggregation than by \textbf{learned} attention weights under this
budget, and that attention adds \textbf{parameters} without compensating
gains. \textbf{Not attempted here:} global transformers, latent graph
learning, edge-conditioned attention, or wide hyperparameter sweeps (see
§6).

\textbf{Why do social edges lower AUC?} The \textbf{directional
ablation} (proximity-only on a 20-bucket evaluation slice) scores
\textasciitilde0.78 AUC versus \textasciitilde0.73 on the 876-bucket
combined+social run. Because the two are evaluated on different horizons
the absolute gap is \textbf{directional, not like-for-like}, but the
same qualitative pattern---proximity-only ahead of
proximity+social---shows up in every smaller-scale ablation we ran.
Three factors likely contribute:

\begin{enumerate}
\def\labelenumi{\arabic{enumi}.}
\tightlist
\item
  \textbf{Noise from registry grouping.} MID social edges link every
  pair of vessels sharing the same three-digit national prefix. In a
  predominantly Chinese-flagged dataset (MID~412), this connects
  thousands of vessels that never physically co-occur, injecting edges
  that the model must learn to discount.
\item
  \textbf{Over-smoothing.} Adding dense social edges increases the
  effective neighborhood size for each GCN layer. When many neighbors
  are structurally similar but behaviorally unrelated, node embeddings
  blur toward a common mean, reducing the model's ability to distinguish
  genuinely co-present pairs.
\item
  \textbf{Objective dilution.} The training loss treats social edges and
  proximity edges identically (both are positive examples). The model
  therefore spends capacity fitting registry structure rather than
  geographic co-occurrence, which is the signal that carries over to
  unseen test days.
\end{enumerate}

Adding dense MID-based social edges was a deliberate \textbf{first-cut}
way to inject registry context, and the modest performance drop is
informative rather than a dead end: it suggests the right answer is not
``more edges of the same kind'' but \textbf{a different way to fuse
identity information with movement}. Several concrete directions look
more promising and are listed under §6:

\begin{enumerate}
\def\labelenumi{\arabic{enumi}.}
\tightlist
\item
  \textbf{Heterogeneous-edge architectures.} Treat proximity and
  registry as \textbf{different relations} so the model can learn
  separate aggregation functions per edge type (R-GCN, HGT, or a typed
  message-passing layer). The current homogeneous GCN forces both
  relations through the same weight matrix, which is exactly where the
  dilution arises.
\item
  \textbf{Edge-type-specific weights inside the loss.} A simple
  intermediate step before changing the architecture: down-weight social
  positives (or sample them less aggressively) so that
  \textbf{proximity} edges dominate the training signal while social
  edges still contribute as a soft prior.
\item
  \textbf{MID as a node feature, not an edge.} Encoding the three-digit
  MID prefix as a learned \textbf{node embedding} concatenated with the
  temporal node features lets the model use registry context as a
  \textbf{bias} on each vessel rather than as a forced positive between
  every same-flag pair. This avoids creating thousands of structurally
  identical neighbors in MID 412 while still letting the encoder
  condition on flag.
\item
  \textbf{Tighter or richer registry signals.} Beyond MID, ownership
  clusters, port-of-registry, or operator IDs (where available) are
  usually more informative than three-digit national prefixes; combined
  with (1) and (3), they would also let the model learn to
  \textbf{discount} very common groupings.
\end{enumerate}

The point is not that registry information is unhelpful---it is that
\textbf{how} it enters the graph matters at least as much as
\textbf{whether} it enters at all.

\subsection{Geographic validation and case
studies}\label{geographic-validation-and-case-studies}

Under the workflow that ranks candidates and then applies a
\textbf{25~km}, \textbf{±1~day} rule, \textbf{eight} pairs remain. Two
examples are \textbf{412422375 \(\leftrightarrow\) 412428225}
(\textbf{2} qualifying days) and \textbf{412000690 \(\leftrightarrow\)
412325200} (\textbf{102} days). Extended discussion and plots appear in
\texttt{docs/candidate\_case\_studies.md}. §4.2.1 summarizes
\textbf{gear composition} across all \textbf{200} ranked pairs and
relates the \textbf{192} excluded pairs to \textbf{expectations} from
fleet-wide co-location.

\textbf{Figure~1 --- purpose.} Link scores are abstract; \textbf{maps}
tie them to geography. \textbf{Figure~1} highlights the pair
\textbf{412000690 \(\leftrightarrow\) 412325200} (\textbf{102} days
within \textbf{25~km}) against the \textbf{Chinese} coast
(\textbf{Yellow Sea / East China Sea}). Chance pairs would rarely show
\textbf{persistent} spatial overlap; these tracks \textbf{do}
overlap---supporting that the score reflects \textbf{real co-movement},
not that we have proven illicit coordination. \textbf{Table~2} and
Figures~3--4 summarize \textbf{eight} pairs under a \textbf{fixed}
full-coverage screening rule (the pair set differs slightly from
Figure~1); Figure~1 remains the clearest \textbf{single-pair}
illustration.

How to read the figure:

\begin{itemize}
\tightlist
\item
  \textbf{Blue lines} trace every daily mean position for vessel
  \textbf{412000690}; \textbf{orange lines} trace vessel
  \textbf{412325200}, both spanning 2013 to 2018.
\item
  The \textbf{colored background} is a \textbf{combined daily presence
  density}: for each vessel a smoothed density surface (2D kernel
  density estimation) is computed from all of its daily mean positions,
  then the two surfaces are \textbf{summed} and normalized to a 0--1
  scale. The horizontal color bar at the bottom maps this scale:

  \begin{itemize}
  \tightlist
  \item
    \textbf{Dark blue / purple} (0.00--0.15) = low combined presence ---
    neither vessel spent much time there.
  \item
    \textbf{Teal / green} (0.15--0.45) = moderate combined presence ---
    at least one vessel was present occasionally.
  \item
    \textbf{Yellow / orange / red} (0.45--1.00) = high combined presence
    --- one or both vessels were frequently located there; the warmest
    colors mark the densest overlap.
  \end{itemize}
\item
  \textbf{Black contour lines} with numeric labels (0.1, 0.3, etc.)
  outline specific density thresholds, making it easy to locate the core
  shared operating area.
\item
  The \textbf{warmest region} (yellow-green, centered near
  \textasciitilde121 E, 37--39 N off the Shandong Peninsula) is where
  daily positions from \textbf{both} vessels overlap most, consistent
  with a \textbf{shared fishing ground} in the Yellow Sea.
\item
  \textbf{Tan shading} = land (10 m Natural Earth coastline); labeled
  cities (\textbf{Shanghai}, \textbf{Qingdao}, \textbf{Jinan},
  \textbf{Nanjing}, \textbf{Yancheng}) provide geographic reference
  along the coast.
\end{itemize}

\pandocbounded{\includegraphics[keepaspectratio,alt={Case study tracks}]{../artifacts/plots/case_study_pairs/pair_412000690_412325200_contour.png}}\\
\emph{Figure 1. Flagship candidate pair (\textbf{412000690}, blue;
\textbf{412325200}, orange), Yellow Sea / Chinese coast, 2013--2018.
Shading: summed kernel density of daily presence (0--1); contours mark
density thresholds (see §4.2 ``How to read''). Confirms sustained
overlap near \textasciitilde121°E, 37--39°N---link score aligns with a
shared operating area, not random pairing. Tracks-only PNG and commands:
Appendix~B.}

\subsubsection{Gear-aware candidate
stratification}\label{gear-aware-candidate-stratification}

The model nominates about \textbf{200} pairs; after the \textbf{25~km /
±1~day} geographic rule, \textbf{eight} remain. The question is what the
other \textbf{192} pairs represent. A \textbf{gear-aware}
post-processing step summarizes their \textbf{fleet context} and offers
a numerical rationale for treating many of them as \textbf{ordinary}
same-fleet overlap.

\textbf{Method.} For \textbf{267} distinct MMSIs in the top-200 list we
infer a \textbf{dominant gear} by sampling \textbf{96} vessel-days (one
per month from 2012--2019) and joining to \textbf{fleet} grids with the
same \textbf{hours-weighted} rule used elsewhere (§4.3.1).
Independently, we estimate how often each \textbf{gear pair}
(trawler--trawler, trawler--gillnet, etc.) appears in the same grid cell
on the same day using \textbf{32} sampled fleet days spanning
2012--2019---about \textbf{seven million} cell-day observations in
total. That baseline answers how common each gear pairing is
\textbf{before} consulting any model score.

Each candidate receives \texttt{adjusted\ =\ TGCN\_score\ ×\ discount}
with
\texttt{discount\ =\ 1\ /\ (1\ +\ 2\ ×\ normalized\_baseline\_rate)}.
Gear combinations that are \textbf{frequent} in the fleet receive a
\textbf{larger} discount; \textbf{rarer} or \textbf{cross-gear}
combinations are discounted \textbf{less}. Unknown gear receives
\textbf{no} change.

\textbf{Key findings (Figure 2):}

\begin{itemize}
\tightlist
\item
  \textbf{75\% of the top-200 candidates} are trawler-dominated: 82
  pairs (41\%) are \textbf{Trawlers + Trawlers} and 67 (34\%) are
  \textbf{Line / generic + Trawlers}. However the two categories tell
  very different stories: \textbf{Trawlers + Trawlers} is actually
  \emph{1.1× below} the fleet baseline --- the model is not
  preferentially elevating same-gear pairs; there are simply many
  trawler pairs in the data. \textbf{Line / generic + Trawlers}, by
  contrast, appears \textbf{10.5× above} the fleet baseline --- the TGCN
  is specifically surfacing these cross-gear encounters far more than
  chance would predict, suggesting genuine signal.
\item
  \textbf{50\% are same-flag} (predominantly CHN + CHN), consistent with
  vessels from the same national fleet operating in the same waters.
\item
  After gear adjustment, cross-gear validated pairs \textbf{held or
  improved} their rank: the \textbf{Set gillnets + Trawlers} pair
  (412461376 \(\leftrightarrow\) 412415321) jumped \textbf{48 ranks} (72
  → 24), and the \textbf{Fixed gear + Trawlers} pair (412422375
  \(\leftrightarrow\) 412428225) dropped only 2 ranks.
\item
  The two \textbf{Trawlers + Trawlers} validated pairs dropped
  \textbf{70+ ranks} each (e.g., 80 → 154), confirming their co-presence
  is explained by shared fishing grounds rather than distinctive
  coordination.
\end{itemize}

\textbf{Why the results look this way.} The top-200 list is dominated by
trawler pairs for a straightforward reason: trawlers are the most common
vessel type in the dataset and they concentrate in the same high-density
fishing grounds (Yellow Sea, East China Sea). From the TGCN's
perspective, two trawlers that share the same seasonal migration look
structurally identical --- they appear in the same grid cells, on the
same days, year after year. The model correctly identifies this overlap
but cannot distinguish ``two boats following the same fish'' from ``two
boats deliberately meeting.'' That is why 80 of the 82 trawler-trawler
candidates \textbf{fail} the strict 25 km / ±1 day check: their
graph-level similarity does not translate to sustained physical
proximity on a given day.

Cross-gear pairs tell the opposite story. A fixed-gear vessel
(stationary nets or traps) and a trawler (moving nets) have different
operating modes, target different species, and rarely share the same
grid cell by chance. When the model flags such a pair \emph{and} the
geographic filter confirms they were repeatedly within 25 km, the
co-location is unlikely to be coincidental. This is reflected in the
pass rates: \textbf{Fixed gear + Trawlers} passes at \textbf{25\%} (1 of
4), compared to \textbf{2.4\%} for trawler-trawler pairs --- roughly a
\textbf{10× difference}. The intermediate categories follow the same
gradient: \textbf{Set gillnets + Trawlers} passes at \textbf{8.3\%}, and
\textbf{Line / generic + Trawlers} at \textbf{6.0\%}.

\textbf{What this means for the project.} The gear-aware analysis serves
two purposes:

\begin{enumerate}
\def\labelenumi{\arabic{enumi}.}
\tightlist
\item
  \textbf{It justifies the geographic filter.} The 192 excluded pairs
  are not ``false positives'' in the traditional sense --- the model
  genuinely detected shared movement patterns. But those patterns are
  explained by normal fleet behavior (thousands of same-type vessels
  fishing the same waters), not by vessel-to-vessel coordination. The
  filter correctly separates fleet-level noise from pair-specific
  signal.
\item
  \textbf{It strengthens the 8 validated pairs.} The fact that
  cross-gear pairs survive at disproportionately high rates means the
  validated set is enriched for behaviorally unusual encounters ---
  exactly the kind of co-location that warrants further investigation.
  In particular, the Set gillnets + Trawlers pair that jumped 48 ranks
  after adjustment (72 → 24) was originally buried among dozens of
  trawler pairs; the gear-aware lens reveals it as one of the most
  distinctive candidates in the entire list.
\end{enumerate}

How to read Figure 2:

\begin{itemize}
\tightlist
\item
  \textbf{Panel (a)} answers ``why did 192 of 200 fail?'' Each
  horizontal bar is one gear-type combination; bar length is the
  \textbf{number of candidate pairs} in that category (all 200 are
  shown). Each bar is \textbf{stacked}: gray = did not pass the 25 km /
  ±1 day geographic check; green = passed. \textbf{Failed} and
  \textbf{passed} counts are written once to the \textbf{right of each
  bar} (not inside narrow segments), so labels do not overlap. A strip
  below the chart shows the \textbf{whole cohort} (192 failed + 8 passed
  = 200). The bracket still highlights that \textbf{pass rates} (passed
  ÷ category size) are much higher for cross-gear categories than for
  \textbf{Trawlers + Trawlers}---e.g.~\textbf{Fixed gear + Trawlers}
  \textbf{25\%}, \textbf{Set gillnets + Trawlers} \textbf{8.3\%},
  \textbf{Line / generic + Trawlers} \textbf{6.0\%}, versus
  \textbf{2.4\%} for \textbf{Trawlers + Trawlers} (82 pairs), roughly
  \textbf{10×} at the extremes.
\item
  \textbf{Panel (b)} answers ``should we trust the 8 that survived?'' A
  table lists every validated pair with its MMSI identifiers, gear
  combination, flag match, days within 25 km, mean distance, and TGCN
  rank. Six of eight are cross-gear pairs, consistent with genuine
  vessel-to-vessel encounters rather than fleet noise. The two
  trawler-trawler pairs have higher mean distances (191--349 km across
  overlap days) and lower ranks, further supporting the interpretation
  that their validation is marginal compared to the cross-gear pairs.
\end{itemize}

\pandocbounded{\includegraphics[keepaspectratio,alt={Gear-aware re-ranking}]{../artifacts/plots/gear_aware_reranking.png}}\\
\emph{Figure 2. Gear-stratified review of top-200 candidates after the
\textbf{25~km / ±1~day} rule. \textbf{(a)} Counts by gear-type
combination: gray = failed geography, green = passed (all 200 pairs).
\textbf{(b)} Summary table of the eight passing pairs (gear, flag, days
within 25~km, rank). Script: \texttt{scripts/gear\_aware\_reranking.py};
Appendix~B.}

\subsection{Monthly overlap and cluster
context}\label{monthly-overlap-and-cluster-context}

\textbf{Figure 3} charts the monthly count of days each validated pair's
vessels were within \textbf{25 km}. \textbf{Figure 4} maps those same
rendezvous locations geographically and highlights a \textbf{six-vessel}
subgroup near \textbf{\textasciitilde30 N, 122 E}.

How to read Figure 3:

\begin{itemize}
\tightlist
\item
  Each \textbf{row} is one vessel pair, sorted from most overlap (top)
  to least (bottom). The \textbf{y-axis label} shows the pair number,
  the flag state(s) involved (e.g., CHN for China), and a summary line
  giving the total days and active months.
\item
  Each \textbf{bubble} represents one calendar month. \textbf{Bubble
  size is proportional to the number of days} the two vessels were
  within 25 km that month (see the size legend in the upper right: 1, 5,
  10, 20 days).
\item
  \textbf{Numbers inside larger bubbles} give the exact day count for
  that month.
\item
  \textbf{Blank spaces} mean zero close approaches that month --- the
  pair was either not active or never came within 25 km.
\item
  \textbf{Alternating row shading} (white / light gray) helps visually
  separate pairs.
\item
  \textbf{Color} distinguishes pairs but carries no additional meaning
  beyond identification.
\end{itemize}

How to read Figure 4:

\begin{itemize}
\tightlist
\item
  \textbf{Faint gray dots} show all rendezvous points (monthly locations
  where any of the eight pairs were within 25 km).
\item
  \textbf{Colored markers} (circles, squares, diamonds, triangles,
  crosses) each represent one of the \textbf{six cluster vessels}. The
  legend in the upper left identifies each vessel by abbreviated MMSI
  and flag state (all CHN).
\item
  \textbf{Thin gray lines} connect cluster members that appear together
  as a pair, showing which vessels rendezvous with which.
\item
  The \textbf{red X} marks the \textbf{weighted hotspot centroid} ---
  the average location of all cluster rendezvous points
  (\textasciitilde29.7 N, 122.3 E).
\item
  \textbf{Tan land shading}, coastlines, rivers, and labeled cities
  (\textbf{Shanghai}, \textbf{Hangzhou}, \textbf{Ningbo},
  \textbf{Wenzhou}, \textbf{Fuzhou}) provide geographic reference.
\end{itemize}

\subsubsection{Gear type codes (how to read ``country ·
gear'')}\label{gear-type-codes-how-to-read-country-gear}

Figures~3--4 and Tables~2--3 label each vessel as \textbf{country ·
gear}. \textbf{Country} is taken from the enrichment file
\texttt{artifacts/cooperative\_pairs\_with\_flag\_gear.csv} when
present; otherwise it is inferred from the \textbf{ITU MID} prefix (for
example \textbf{MID111}). \textbf{Gear} strings such as
\texttt{trawlers} or \texttt{set\_gillnets} come from \textbf{fleet}
\texttt{geartype} fields on a \textbf{grid} and are assigned by
\textbf{hours-weighted} overlap between vessel tracks and fleet cells
(\texttt{scripts/enrich\_pairs\_with\_flag\_gear.py}). These labels
describe \textbf{coarse fisheries-activity categories}, \textbf{not}
formal licensing or gear certificates. Tables~2--3 and the figures use
the \textbf{2012--2019} enrichment---the same basis as §4.2.1---which is
\textbf{broader} than an August~2017-only enrichment. Definitions for
each code appear in \textbf{Appendix~A}.

\textbf{Table 2 --- Eight validated pairs (country · gear per vessel).}
\emph{Format: src MMSI, dst MMSI, then src country · gear \textbar{} dst
country · gear.}

{\def\LTcaptype{none} % do not increment counter
\begin{longtable}[]{@{}lll@{}}
\toprule\noalign{}
src & dst & Country · gear (src \textbar{} dst) \\
\midrule\noalign{}
\endhead
\bottomrule\noalign{}
\endlastfoot
412422375 & 412428225 & CHN · fixed\_gear \textbar{} CHN · trawlers \\
412461376 & 412415321 & CHN · trawlers \textbar{} CHN · set\_gillnets \\
412437423 & 412435485 & CHN · trawlers \textbar{} CHN · fishing \\
412410128 & 412416248 & CHN · fishing \textbar{} CHN · trawlers \\
412420679 & 412413383 & CHN · fishing \textbar{} CHN · trawlers \\
412450427 & 111203412 & CHN · fishing \textbar{} MID111 · trawlers \\
412985698 & 412443375 & CHN · trawlers \textbar{} CHN · trawlers \\
412461376 & 412427825 & CHN · trawlers \textbar{} CHN · trawlers \\
\end{longtable}
}

\textbf{Table 3 --- Six-vessel cluster (Figure 4).} \emph{MMSIs selected
by \texttt{scripts/analyze\_six\_vessel\_cluster.py}; labels match plot
legend.}

{\def\LTcaptype{none} % do not increment counter
\begin{longtable}[]{@{}ll@{}}
\toprule\noalign{}
MMSI & Country · gear \\
\midrule\noalign{}
\endhead
\bottomrule\noalign{}
\endlastfoot
412413383 & CHN · trawlers \\
412420679 & CHN · fishing \\
412427825 & CHN · trawlers \\
412435485 & CHN · fishing \\
412437423 & CHN · trawlers \\
412461376 & CHN · trawlers \\
\end{longtable}
}

\pandocbounded{\includegraphics[keepaspectratio,alt={All pairs days within 25 km (monthly)}]{../artifacts/plots/pair_overlap_series/all_pairs_days_within_25km.png}}\\
\emph{Figure 3. Monthly days within \textbf{25~km} for the eight
validated pairs (2013--2019). One row per pair (sorted by total
overlap); bubble area is proportional to days that month (see legend).
Axis labels give total overlap days and months active. Pair~1 shows
dense 2015--2016 overlap; others are mostly sporadic; pair~7 is the only
cross-flag case. Reading guide: §4.3.}

\pandocbounded{\includegraphics[keepaspectratio,alt={Six-vessel cluster}]{../artifacts/plots/six_vessel_cluster_scatter.png}}\\
\emph{Figure 4. Rendezvous locations (monthly, within 25~km) for
validated pairs, with Chinese coast context (Fuzhou--Shanghai). Gray
dots: all rendezvous points; colored symbols: six vessels that recur
across pairs (legend abbreviates MMSI); gray segments connect recorded
pairings. Red \textbf{×}: centroid of cluster rendezvous
(\textasciitilde29.7°N, 122.3°E), east of Ningbo. Highlights repeated
convergence of multiple pairs in one offshore band---consistent with
Figure~3.}

\subsection{Exploratory short-window analysis (August 2017; not the main
benchmark)}\label{exploratory-short-window-analysis-august-2017-not-the-main-benchmark}

Early in the project we trained a small TGCN on \textbf{five}
consecutive days in \textbf{August 2017} (7--11 Aug)---one of the few
intervals where cell-level fleet products provided \textbf{country and
gear} labels for enrichment. Metrics on that slice (about \textbf{0.56}
AUC, \textbf{0.67} AP) sit \textbf{far below} the 2012--2019 benchmark
and should be treated only as \textbf{illustrative}.

\textbf{Why we abandoned the August~2017 shortcut for the main
pipeline.} The five-day window was attractive at first because the
corresponding fleet enrichment file came pre-labeled, which made gear
and country joins trivial; but five days of test data is \textbf{too
short} for a temporal benchmark (only \textbf{5} test buckets versus
\textbf{876} in the headline run), concentrates on a single seasonal
regime that does \textbf{not} reflect annual variability, and ties
\textbf{gear assignment} to whatever vessels happened to be active in
that one summer week. The full-timeline gear analysis in §4.2.1
\textbf{replaces} this shortcut: it assigns gear using
\textbf{2012--2019} fleet context for all \textbf{200}
candidates---roughly \textbf{seven million} cell-day observations across
\textbf{32} sampled fleet days---so the labels reflect each vessel's
\textbf{dominant} behavior over years rather than a single-week
snapshot, and the headline metrics in Table~1 sit on a properly long
evaluation horizon. An August~2017 cooperation graphic remains in
\texttt{artifacts/plots/tgcn\_daily\_cooperation\_heatmap.png} for
archival purposes but is \textbf{not} part of this document. Gear codes
are defined in \textbf{Appendix~A}.

\begin{center}\rule{0.5\linewidth}{0.5pt}\end{center}

\section{Discussion}\label{discussion}

This section states \textbf{limitations} of the task, data, metrics, and
operational use---standard before §6.

\subsection{Threats and limitations}\label{threats-and-limitations}

The same concerns as in earlier drafts, grouped so readers can scan by
\textbf{what kind} of limitation is at issue.

\subsubsection{Task definition and
semantics}\label{task-definition-and-semantics}

\begin{itemize}
\tightlist
\item
  \textbf{Link prediction ranks edges under a graph rule, not
  misconduct.} ROC AUC / AP measure fit to observed vs sampled
  \textbf{non-edges}. High scores follow from \textbf{shared grounds},
  \textbf{lanes}, or \textbf{ports} as easily as from rare
  coordination.\\
\item
  \textbf{No cooperation, IUU, or ``dark fleet'' labels} were used for
  training; the objective is structural co-occurrence only.\\
\item
  \textbf{Shared geography is not the same as coordination.} In crowded
  fisheries, overlap is \textbf{ambiguous} even when it is wholly legal
  and routine.
\end{itemize}

\subsubsection{Data and preprocessing}\label{data-and-preprocessing}

\begin{itemize}
\tightlist
\item
  \textbf{AIS} can be missing, intermittent, or misused; graphs only see
  \textbf{broadcast} tracks.\\
\item
  \textbf{Daily means, grid cells, and merge rules} shape both edges and
  post-hoc distances.\\
\item
  \textbf{Dense regions} yield many possible edges---\textbf{degree} and
  \textbf{baseline encounter rate} can inflate scores relative to
  detecting \textbf{rare} pairwise events.\\
\item
  The \textbf{August~2017} slice is a \textbf{narrow} exploratory
  window---\textbf{do not} treat it like the 2012--2019 benchmark.\\
\item
  \textbf{Country · gear} come from \textbf{fleet cell enrichment}
  (§4.3.1, \textbf{Appendix~A}): useful for narrative context,
  \textbf{not} a substitute for registry or license records.
\end{itemize}

\subsubsection{Evaluation and
reproducibility}\label{evaluation-and-reproducibility}

\begin{itemize}
\tightlist
\item
  \textbf{Metrics are not enforcement decisions.} Good ranking on the
  surrogate task does \textbf{not} validate real-world outcomes;
  \textbf{ground-truth coordination} would be needed and is
  \textbf{absent} here.\\
\item
  The headline combined-graph row is \textbf{single-seed}; §4.1 gives
  day-level spread and seed spread on \textbf{other}
  setups---\textbf{multi-seed} uncertainty on the full benchmark is
  \textbf{future work}.\\
\item
  \textbf{Gear-aware adjustment} (§4.2.1) uses a \textbf{sampled} fleet
  baseline (\textbf{32} fleet days, \textasciitilde7~million cell-days)
  and a \textbf{fixed discount} formula: \textbf{transparent},
  \textbf{not} calibrated; alternate formulas would \textbf{change
  ranks}. Use it to \textbf{interpret} the long list, \textbf{not} as a
  second detector replacing geography.
\end{itemize}

\subsubsection{Deployment and computational
constraints}\label{deployment-and-computational-constraints}

\begin{itemize}
\tightlist
\item
  \textbf{Training} was limited to \textbf{\textasciitilde1450} days on
  a \textbf{16~GB} laptop; larger caps hit \textbf{memory}
  failure---reported limits are \textbf{hardware-bound}, not an
  intrinsic ceiling on the method.\\
\item
  In operations, treat outputs as \textbf{ranked lists} plus
  maps/tables: \textbf{analyst triage}, \textbf{not} automated findings.
\end{itemize}

\subsection{Synthesis}\label{synthesis}

\textbf{Social edges} link vessels by registry prefix even when they
rarely meet; that injects \textbf{noise}, encourages
\textbf{over-smoothing}, and \textbf{dilutes} proximity in the training
loss---so the \textbf{combined} graph's headline scores \textbf{trail}
the \textbf{proximity-only} curve despite embedding organizational
signal one might encode differently in another architecture.
\textbf{Strict post-hoc geography} (25~km, ±1~day on daily means) tests
whether pairs flagged by \textbf{10~km} same-day structure actually
\textbf{repeatedly co-locate} at operational distances; it
\textbf{separates} fleet-scale \textbf{structural similarity} from
\textbf{pair-level} proximity. \textbf{``Validated''} here means
\textbf{passing that screen} together with documented \textbf{tracks and
overlap}---an analyst gets a \textbf{short} list of \textbf{reviewable}
cases with \textbf{reproducible exhibits}, \textbf{not} a determination
of guilt or IUU.

\begin{center}\rule{0.5\linewidth}{0.5pt}\end{center}

\section{Conclusion and future work}\label{conclusion-and-future-work}

This report describes an \textbf{unsupervised} pipeline from AIS to
daily graphs to a \textbf{TGCN} for \textbf{temporal link prediction},
evaluated with a \textbf{chronological} train/test split (§3.1; formal
symbols in \textbf{Appendix~C}). On the \textbf{combined}
graph---proximity edges plus \textbf{MID-consistent} social edges---the
documented laptop run reaches \textbf{ROC AUC = 0.72997} and \textbf{AP
= 0.71381} on \textbf{876} test days (\textbf{Table~1}), with
\textbf{proximity-only} and \textbf{capped-graph} settings reported for
comparison. \textbf{Eight} pairs remain after conservative
\textbf{post-hoc} distance filters (\textbf{Table~2}). The
\textbf{gear-aware} analysis (§4.2.1) suggests that much of the top-200
list reflects \textbf{routine} trawler-dominated fleet structure,
whereas pairs that clear geography---and especially \textbf{cross-gear}
pairs---look \textbf{less} typical relative to fleet baselines. Figures
and tables add \textbf{country and gear} context where enrichment exists
(\textbf{Appendix~A}). For operational use, the deliverable is a
\textbf{ranked candidate list} with \textbf{reproducible} geographic and
contextual checks---not a finding of guilt or confirmed IUU.

\textbf{Future work.} Three directions stand out as both feasible on the
existing pipeline and high-value:

\begin{enumerate}
\def\labelenumi{\arabic{enumi}.}
\tightlist
\item
  \textbf{Multi-seed uncertainty on the full benchmark.} Re-run the
  headline combined-graph configuration with at least \textbf{3--5
  seeds} so the \textbf{single-seed} headline of \textbf{0.72997 /
  0.71381} can be reported with a proper confidence interval rather than
  only the day-level spread (§4.1). Pair this with a
  \textbf{precision-at-K} sweep (top \textbf{50 / 100 / 200} pairs
  versus geographic survival rate) so the \textbf{operational} value of
  the model is reported in the same table as the macro metrics.
\item
  \textbf{Better fusion of registry / ownership information.} The
  observed drop from the proximity-only ablation
  (\textasciitilde{}\textbf{0.78} AUC on a 20-bucket slice) to the
  combined+social run (\textasciitilde{}\textbf{0.73} AUC on 876
  buckets) does not say that registry context is unhelpful---it says
  that \textbf{dense, equally-weighted} social edges between every
  same-MID pair are the wrong way to inject it. More promising fusions,
  all listed in §4.1: \textbf{heterogeneous-edge} architectures (R-GCN /
  HGT / typed message passing) that learn a \textbf{separate}
  aggregation per edge type; \textbf{edge-type-specific weights} in the
  loss so proximity dominates the training signal; and treating MID (or
  richer ownership / port-of-registry IDs) as a \textbf{node feature}
  concatenated with the temporal embedding rather than as an explicit
  edge---conditioning the encoder on flag without creating thousands of
  structurally identical neighbors. Each of these should also tighten
  the \textbf{top-200 → geography → gear-context} funnel in
  \textbf{Table 1b}, raising the analyst-yield rate above the current
  \textbf{4\%} and---if the gap between trawler--trawler and
  best-cross-gear pass rates persists---make the operational top-N list
  cleaner without an analyst having to filter post-hoc.
\item
  \textbf{Scaling and richer signals.} Finer spatial-temporal
  resolution; \textbf{supervised or semi-supervised} extensions where
  trustworthy labels exist (e.g., declared transshipment events);
  \textbf{movement-based} gear or behavior models; larger day budgets on
  \textbf{high-memory} machines beyond the
  \textbf{\textasciitilde1450}-day laptop ceiling; and a disciplined
  separation of the \textbf{main 2012--2019 benchmark} from
  \textbf{illustrative} short windows such as August 2017.
\end{enumerate}

\clearpage

\section*{References}\label{references}
\addcontentsline{toc}{section}{References}

\begin{itemize}
\tightlist
\item
  Kroodsma, D. A., et al.~(2018). \emph{Tracking the global footprint of
  fisheries}. Science, 359(6378), 904--908.\\
\item
  Miller, N. A., et al.~(2018). \emph{Identifying global patterns of
  transshipment behavior}. Frontiers in Marine Science, 5, 240.\\
\item
  Park, J., et al.~(2020). \emph{A systematic assessment of vessel
  monitoring data for identifying suspicious transshipment events}. ICES
  Journal of Marine Science.\\
\item
  Kipf, T. N., \& Welling, M. (2017). \emph{Semi-Supervised
  Classification with Graph Convolutional Networks}. ICLR.\\
\item
  Hamilton, W. L., Ying, R., \& Leskovec, J. (2017). \emph{Inductive
  Representation Learning on Large Graphs}. NeurIPS.\\
\item
  Yu, B., Yin, H., \& Zhu, Z. (2018). \emph{Spatio-Temporal Graph
  Convolutional Networks: A Deep Learning Framework for Traffic
  Forecasting}. IJCAI.\\
\item
  Seo, Y., et al.~(2018). \emph{Structured Sequence Modeling with Graph
  Convolutional Recurrent Networks}. ICONIP.\\
\item
  Rossi, E., et al.~(2020). \emph{Temporal Graph Networks for Deep
  Learning on Dynamic Graphs}. ICML Workshop on Graph Representation
  Learning.\\
\item
  Lü, L., \& Zhou, T. (2011). \emph{Link prediction in complex networks:
  A survey}. Physica A, 390(6), 1150--1170.
\end{itemize}

\begin{center}\rule{0.5\linewidth}{0.5pt}\end{center}

\clearpage

\section*{Appendix A --- Gear type
codes}\label{appendix-a-gear-type-codes}
\addcontentsline{toc}{section}{Appendix A --- Gear type codes}

\emph{Plain-language definitions for \textbf{gear} strings in
\textbf{Tables 2--3} and \textbf{Figures 2--4}. The same content is kept
in the repository as \textbf{\texttt{docs/gear\_types.md}} for version
control and easy editing.}

\subsection*{Source of labels}\label{source-of-labels}
\addcontentsline{toc}{subsection}{Source of labels}

Figure and table \textbf{gear} labels are \textbf{not} copied from a
vessel registry. For each vessel--day we intersect AIS grid cells with
\textbf{fleet} rows that report \textbf{hours} by \textbf{flag ×
geartype} within the cell; we assign the \textbf{geartype with the
largest weighted hours}
(\texttt{scripts/enrich\_pairs\_with\_flag\_gear.py}). The result is a
\textbf{coarse behavioral tag}---``what fishing activity dominates that
cell's attributed effort''---rather than an official gear permit for
that MMSI.

Plots show \textbf{---} when enrichment fails (no fleet overlap or
vessel absent from the join).

\subsection*{\texorpdfstring{Codes in
\texttt{artifacts/cooperative\_pairs\_with\_flag\_gear.csv}}{Codes in artifacts/cooperative\_pairs\_with\_flag\_gear.csv}}\label{codes-in-artifactscooperative_pairs_with_flag_gear.csv}
\addcontentsline{toc}{subsection}{Codes in
\texttt{artifacts/cooperative\_pairs\_with\_flag\_gear.csv}}

These strings are taken \textbf{as-is} from the fleet
\textbf{\texttt{geartype}} column (snake\_case).

{\def\LTcaptype{none} % do not increment counter
\begin{longtable}[]{@{}
  >{\raggedright\arraybackslash}p{(\linewidth - 2\tabcolsep) * \real{0.4800}}
  >{\raggedright\arraybackslash}p{(\linewidth - 2\tabcolsep) * \real{0.5200}}@{}}
\toprule\noalign{}
\begin{minipage}[b]{\linewidth}\raggedright
Code (CSV)
\end{minipage} & \begin{minipage}[b]{\linewidth}\raggedright
Short gloss
\end{minipage} \\
\midrule\noalign{}
\endhead
\bottomrule\noalign{}
\endlastfoot
\textbf{\texttt{fishing}} & \textbf{Generic / unspecified fishing.}
Catch-all category in the source data when activity is attributed to
``fishing'' without a more specific gear class. Treat as \textbf{low
specificity} compared to trawlers, nets, etc. \\
\textbf{\texttt{trawlers}} & \textbf{Trawling.} Vessels \textbf{tow
nets} through the water column or along the bottom (e.g.~otter trawl,
beam trawl). High mobility; distinct from \textbf{fixed} or \textbf{set}
gear. \\
\textbf{\texttt{fixed\_gear}} & \textbf{Stationary gear} fixed to the
seabed or structure: pots, traps, stakes, weirs, and similar
\textbf{non-towed} gear that stays in place. (Name reflects ``fixed''
position, not ``repaired.'') \\
\textbf{\texttt{set\_gillnets}} & \textbf{Set gillnets.} Nets
\textbf{set and left} (anchored, on bottom, or sometimes drifting) so
fish \textbf{gill} in the mesh. Not actively towed like a trawl. \\
\textbf{\texttt{set\_longlines}} & \textbf{Set longlines.} A
\textbf{long line} with many \textbf{baited hooks}, deployed and
\textbf{left to soak} (demersal or pelagic longline, depending on
fishery). \\
\textbf{\texttt{other\_purse\_seines}} & \textbf{Purse seining (other).}
A \textbf{surrounding net} used on \textbf{schooling fish}; the bottom
is \textbf{pulled closed} (``pursed'') like a drawstring. ``Other''
indicates a sub-type bucket in the source taxonomy (not necessarily
``miscellaneous quality''). \\
\textbf{\texttt{pole\_and\_line}} & \textbf{Pole-and-line / baitboat.}
Fish caught with \textbf{hand-held poles} and \textbf{hooks}, often with
\textbf{live bait}; common in some tuna fisheries. \\
\end{longtable}
}

\subsection*{How to interpret labels}\label{how-to-interpret-labels}
\addcontentsline{toc}{subsection}{How to interpret labels}

\begin{enumerate}
\def\labelenumi{\arabic{enumi}.}
\tightlist
\item
  \textbf{Same MMSI, different days} can get different geartypes if the
  vessel moves between cells dominated by different fleet
  attributions.\\
\item
  \textbf{Same geartype on two MMSIs} does not prove they use identical
  gear---only that both were attributed the same \textbf{coarse} class
  in the fleet enrichment window (2012--2019 for all tables and
  figures).\\
\item
  \textbf{AIS ``ship type''} codes differ from fleet \textbf{geartype}
  here: the latter is \textbf{fisheries-activity} oriented, not the full
  IMO/AIS ship-type list.
\end{enumerate}

Methodology of the join and weighting:
\texttt{scripts/enrich\_pairs\_with\_flag\_gear.py} and §4.3--4.4 of
this report.

\subsection*{Related work (not from the same
CSV)}\label{related-work-not-from-the-same-csv}
\addcontentsline{toc}{subsection}{Related work (not from the same CSV)}

The repository also includes a \textbf{movement-based gear classifier}
trained on \textbf{anonymized} labeled tracks (see project
\texttt{README.md}, ``Gear classification''). Those class labels
(e.g.~\texttt{purse\_seines}, \texttt{trollers}) are \textbf{not}
automatically the same strings as the fleet \textbf{\texttt{geartype}}
column above; do not merge them without an explicit mapping.

\begin{center}\rule{0.5\linewidth}{0.5pt}\end{center}

\clearpage

\section*{Appendix B ---
Reproducibility}\label{appendix-b-reproducibility}
\addcontentsline{toc}{section}{Appendix B --- Reproducibility}

{\def\LTcaptype{none} % do not increment counter
\begin{longtable}[]{@{}
  >{\raggedright\arraybackslash}p{(\linewidth - 2\tabcolsep) * \real{0.2727}}
  >{\raggedright\arraybackslash}p{(\linewidth - 2\tabcolsep) * \real{0.7273}}@{}}
\toprule\noalign{}
\begin{minipage}[b]{\linewidth}\raggedright
Item
\end{minipage} & \begin{minipage}[b]{\linewidth}\raggedright
Command / path
\end{minipage} \\
\midrule\noalign{}
\endhead
\bottomrule\noalign{}
\endlastfoot
Combined edges &
\texttt{python3\ scripts/add\_social\_edges.py\ -\/-edges\ artifacts/edges\_2012\_2019\_full.parquet\ -\/-out\ artifacts/edges\_full\_with\_social.parquet\ -\/-max-social-per-bucket\ 2000} \\
Primary TGCN JSON/CSV &
\texttt{artifacts/tgcn\_social\_maxb1450\_ep3.json}, \texttt{.csv} ---
PyG env, \texttt{PYTHONPATH=scripts} \\
Optional per-bucket metric plots (not in main report) &
\texttt{python3\ scripts/plot\_tgcn\_bucket\_metrics.py\ -\/-report\ artifacts/tgcn\_social\_maxb1450\_ep3.json\ -\/-out-dir\ artifacts/plots} \\
Case study tracks (Figure 1) --- command &
\texttt{python3\ scripts/compute\_pair\_overlap\_from\_daily.py\ -\/-pairs\ artifacts/case\_study\_pairs.csv\ -\/-daily-root\ "data/MMSI\ daily\ vessels\ "\ -\/-top-k\ 4\ -\/-distance-km\ 25\ -\/-day-window\ 1\ -\/-max-files-per-year\ 0\ -\/-out-dir\ artifacts/plots/case\_study\_pairs\ -\/-contour} \\
Case study tracks --- notes &
\textbf{\texttt{-\/-max-files-per-year\ 0}} uses all daily CSVs
(default); smaller values can undercount vs
\texttt{docs/candidate\_case\_studies.md}.
\textbf{\texttt{-\/-contour}}: KDE density map (\texttt{Spectral\_r}).
With \textbf{cartopy}: coastlines / land / cities. Outputs
\texttt{pair\_\textless{}src\textgreater{}\_\textless{}dst\textgreater{}\_contour.png}
(+ tracks-only \texttt{.png}). Optional \texttt{-\/-out-summary},
\texttt{-\/-contour-all-pairs}. \\
Ablations & \texttt{scripts/run\_tgcn\_improvement\_suite.py} →
\texttt{artifacts/tgcn\_improvement\_suite\_summary.csv} \\
August cooperation summary (not in report; see §4.4) &
\texttt{PYTHONPATH=scripts\ python3\ scripts/plot\_cooperative\_heatmap.py}
--- two-panel aggregate cooperation summary: stacked bar chart by gear
group + flag-state breakdown (Aug 2017 short window). Output:
\texttt{artifacts/plots/tgcn\_daily\_cooperation\_heatmap.png}. \\
Eight-pair overlap CSV + optional heatmap PNG & Full command in
\textbf{block below} (overlap CSV feeds \textbf{Tables 2--3} / Figures
3--4). \\
Monthly overlap bubble chart (Figure 3) &
\texttt{PYTHONPATH=scripts\ python3\ scripts/plot\_pair\_overlap\_time\_series.py\ -\/-overlap-csv\ artifacts/eight\_pairs\_overlap\_by\_month.csv\ -\/-enrichment\ artifacts/cooperative\_pairs\_with\_flag\_gear.csv}
--- bubble chart of monthly close approaches, one row per pair. \\
Six-vessel cluster map (Figure 4) &
\texttt{PYTHONPATH=scripts\ python3\ scripts/analyze\_six\_vessel\_cluster.py\ -\/-overlap-csv\ artifacts/eight\_pairs\_overlap\_by\_month.csv\ -\/-enrichment\ artifacts/cooperative\_pairs\_with\_flag\_gear.csv}
--- geographic scatter with cartopy coastline, land fill, cities. \\
Enrichment CSV (for labels) &
\texttt{artifacts/cooperative\_pairs\_with\_flag\_gear.csv} ---
\texttt{scripts/enrich\_pairs\_with\_flag\_gear.py} \\
\textbf{Gear type definitions (for tables \& figures)} &
\textbf{Appendix A} (mirror for repo edits:
\texttt{docs/gear\_types.md}) \\
\textbf{Notation / formal link-prediction setup} & \textbf{Appendix C}
(symbol table; training vs validation distances) \\
Gear × country pair-count heatmap (updated PNG) &
\texttt{PYTHONPATH=scripts\ python3\ scripts/plot\_flag\_gear\_enrichment.py\ -\/-input\ artifacts/cooperative\_pairs\_with\_flag\_gear.csv\ -\/-out-dir\ artifacts/plots} \\
Gear-aware re-ranking (Figure 2) &
\texttt{PYTHONPATH=scripts\ python3\ scripts/gear\_aware\_reranking.py}
--- enriches top-200 candidate MMSIs with gear type (96 sampled days),
computes fleet co-location baselines (\textasciitilde7M cell-days), and
produces gear-adjusted rankings + combined bar-chart-and-table
figure. \\
Candidates / close pairs &
\texttt{artifacts/tgcn\_candidate\_scores\_fullcoverage.parquet},
\texttt{artifacts/close\_pairs\_fullcoverage\_25km\_w1.csv} \\
\end{longtable}
}

\textbf{Eight-pair overlap heatmap --- copy/paste (project root):}

\begin{Shaded}
\begin{Highlighting}[]
\VariableTok{PYTHONPATH}\OperatorTok{=}\NormalTok{scripts }\ExtensionTok{python3}\NormalTok{ scripts/overlap\_by\_month\_8pairs.py }\DataTypeTok{\textbackslash{}}
  \AttributeTok{{-}{-}pairs}\NormalTok{ artifacts/close\_pairs\_fullcoverage\_25km\_w1.csv }\DataTypeTok{\textbackslash{}}
  \AttributeTok{{-}{-}daily{-}root} \StringTok{"data/MMSI daily vessels "} \DataTypeTok{\textbackslash{}}
  \AttributeTok{{-}{-}all{-}files} \AttributeTok{{-}{-}full{-}months} \DataTypeTok{\textbackslash{}}
  \AttributeTok{{-}{-}distance{-}km}\NormalTok{ 25 }\AttributeTok{{-}{-}day{-}window}\NormalTok{ 1 }\DataTypeTok{\textbackslash{}}
  \AttributeTok{{-}{-}out{-}csv}\NormalTok{ artifacts/eight\_pairs\_overlap\_by\_month.csv }\DataTypeTok{\textbackslash{}}
  \AttributeTok{{-}{-}out{-}plot}\NormalTok{ artifacts/plots/eight\_pairs\_overlap\_by\_month.png }\DataTypeTok{\textbackslash{}}
  \AttributeTok{{-}{-}enrichment}\NormalTok{ artifacts/cooperative\_pairs\_with\_flag\_gear.csv}
\end{Highlighting}
\end{Shaded}

\emph{Faster test run (subsamples daily files; shorter runtime):} omit
\texttt{-\/-all-files} and add e.g.~\texttt{-\/-max-files-per-year\ 30}.

\begin{center}\rule{0.5\linewidth}{0.5pt}\end{center}

\clearpage

\section*{Appendix C --- Notation and formal problem
setup}\label{appendix-c-notation-and-formal-problem-setup}
\addcontentsline{toc}{section}{Appendix C --- Notation and formal
problem setup}

\emph{This appendix condenses §3.1 for readers who want a symbol list.
It may be \textbf{omitted in slide versions} or reordered after A--B if
a thesis program requires a specific appendix sequence.}

\subsection*{Symbol glossary}\label{symbol-glossary}
\addcontentsline{toc}{subsection}{Symbol glossary}

{\def\LTcaptype{none} % do not increment counter
\begin{longtable}[]{@{}
  >{\raggedright\arraybackslash}p{(\linewidth - 2\tabcolsep) * \real{0.4706}}
  >{\raggedright\arraybackslash}p{(\linewidth - 2\tabcolsep) * \real{0.5294}}@{}}
\toprule\noalign{}
\begin{minipage}[b]{\linewidth}\raggedright
Symbol
\end{minipage} & \begin{minipage}[b]{\linewidth}\raggedright
Meaning
\end{minipage} \\
\midrule\noalign{}
\endhead
\bottomrule\noalign{}
\endlastfoot
\(t\) & A calendar \textbf{day} (time bucket) in the AIS timeline. \\
\(\mathcal{T}\) & The set of days used after filtering (e.g., days with
edges). \\
\(V\) & Set of \textbf{nodes} (9-digit vessel \textbf{MMSIs} appearing
in the extract). \\
\(G_t = (V, E_t)\) & \textbf{Undirected} snapshot graph for day
\(t\). \\
\(E_t\) & \textbf{Edges} for day \(t\): proximity (§3.3) ± optional
social edges (§3.3). \\
\(u, v\) & Distinct nodes (unordered pair \(\{u,v\}\)). \\
\(s_{u,v}\) & \textbf{Scalar score} (logit) for pair \((u,v)\) on a
given forward pass / bucket. \\
\(\mathbf{h}_u\) & \textbf{Embedding vector} for node \(u\) after the
TGCN (dimension \textbf{32} in the primary run). \\
\(\mathcal{T}_{\mathrm{train}}, \mathcal{T}_{\mathrm{test}}\) &
\textbf{Chronological} partition of buckets (\textbf{\textasciitilde70\%
/ \textasciitilde30\%}). \\
\end{longtable}
}

\subsection*{Graph semantics}\label{graph-semantics}
\addcontentsline{toc}{subsection}{Graph semantics}

For each day \(t \in \mathcal{T}\), the graph \(G_t = (V, E_t)\) is
\textbf{undirected}. An edge \(\{u,v\} \in E_t\) means the vessels met
the \textbf{construction rule} for that day---typically
\textbf{same-day} positions within \textbf{10~km}, optionally plus
\textbf{same three-digit MID} ties. Edges encode \textbf{graph
structure}, not adjudicated cooperation. The vertex set \(V\) may be
fixed globally or built per implementation; \(|E_t|\) is kept
\textbf{sparse} and sometimes \textbf{capped per day} for memory (§3.3).

\subsection*{Temporal link prediction objective
(informal)}\label{temporal-link-prediction-objective-informal}
\addcontentsline{toc}{subsection}{Temporal link prediction objective
(informal)}

The model reads the sequence \(\{G_t\}\) (with optional \textbf{temporal
node features}) into recurrent states and node embeddings
\(\mathbf{h}_u\). On each training day in
\(\mathcal{T}_{\mathrm{train}}\), \textbf{positives} are edges present
that day; \textbf{negatives} are sampled non-edges with counts matched
\textbf{per day} in code. Pair logits follow
\(s_{u,v} \propto \mathbf{h}_u^\top \mathbf{h}_v\) and are trained with
\textbf{binary cross-entropy with logits}. Labels encode
\textbf{presence or absence of an edge}, not IUU or transshipment.

\textbf{Evaluation.} On \(\mathcal{T}_{\mathrm{test}}\), ROC AUC and AP
measure how highly \textbf{true} test edges rank versus
\textbf{non-edges} under the evaluation script's protocol. Large values
indicate \textbf{consistent ranking of observed contacts}, not legal
liability.

\subsection*{Training vs validation
distances}\label{training-vs-validation-distances}
\addcontentsline{toc}{subsection}{Training vs validation distances}

{\def\LTcaptype{none} % do not increment counter
\begin{longtable}[]{@{}
  >{\raggedright\arraybackslash}p{(\linewidth - 4\tabcolsep) * \real{0.1750}}
  >{\raggedright\arraybackslash}p{(\linewidth - 4\tabcolsep) * \real{0.6750}}
  >{\raggedright\arraybackslash}p{(\linewidth - 4\tabcolsep) * \real{0.1500}}@{}}
\toprule\noalign{}
\begin{minipage}[b]{\linewidth}\raggedright
Stage
\end{minipage} & \begin{minipage}[b]{\linewidth}\raggedright
Typical distance / rule
\end{minipage} & \begin{minipage}[b]{\linewidth}\raggedright
Role
\end{minipage} \\
\midrule\noalign{}
\endhead
\bottomrule\noalign{}
\endlastfoot
\textbf{Graph edges (training signal)} & \textbf{\(\leq\) 10 km}, same
\textbf{day} & Defines \(E_t\) for learning. \\
\textbf{Candidate screening (post-hoc)} & \textbf{25 / 50 / 100 km},
optional \textbf{±1 day} & \textbf{Not} used as training labels;
validates high-scoring pairs on raw tracks (§3.5). \\
\end{longtable}
}

Sections §3.3--3.5 repeat this distinction so \textbf{training edges}
are not confused with \textbf{analyst thresholds} used afterward.

\end{document}
